\section{Research design}\label{sec:design}

\subsection{Synthetic repeated-sampling benchmark}

The simulation study first asks a question that cannot be answered cleanly with market prices alone: when the data-generating volatility is known, how much is lost by collapsing posterior uncertainty before pricing? Historical return samples are generated under $\mathbb P$ and option values are evaluated under $\mathbb Q$, preserving the measure separation in Section~\ref{sec:methodology}. The original production grid varies historical sample size, volatility, moneyness, and maturity and evaluates posterior-integrated, posterior-mean, posterior-mode, and MLE pricing against the external target
\begin{equation}
 C_0=C^Q(\sigma_0).
 \label{eq:synthetic_target}
\end{equation}
The target is not a posterior draw and is not constructed from any estimator being compared.

For method $M\in\{PI,PM,Mode,MLE\}$ and $R$ independent histories, the main statistics are
\begin{align}
 \operatorname{Bias}_M
 &=R^{-1}\sum_{j=1}^{R}(\widehat C_{M,j}-C_0),\\
 \operatorname{MAE}_M
 &=R^{-1}\sum_{j=1}^{R}|\widehat C_{M,j}-C_0|,\\
 \operatorname{RMSE}_M
 &=\left[R^{-1}\sum_{j=1}^{R}(\widehat C_{M,j}-C_0)^2\right]^{1/2}.
 \label{eq:metrics}
\end{align}

\subsection{Massive mechanism map}

The large mechanism experiment is designed around Equation~\eqref{eq:taylor_gap}. It uses
\begin{align}
 n&\in\{21,42,63,126,252,504,1260\},\\
 \sigma_0&\in\{0.10,0.20,0.35,0.50,0.80\},
\end{align}
with 5,000 independently generated histories in each $(n,\sigma_0)$ cell, for 175,000 synthetic datasets. The contract layer crosses four maturities (21, 63, 126, and 252 days), seven moneyness values from 0.60 to 1.50, and six partial-fixing states from zero to nearly the full averaging window.

Because the Gaussian prior on $\mu$ is conjugate to the likelihood conditional on $\sigma$, $\mu$ is integrated analytically. The remaining posterior in $\sigma$ is evaluated on a dense deterministic grid. This removes random-walk Metropolis error from the mechanism map while retaining exactly the baseline Gaussian statistical model. For each history and contract state the experiment computes $\Delta_{PI,PM}$, posterior variance of $\sigma$, local price curvature, and the Taylor approximation in Equation~\eqref{eq:taylor_gap}. The purpose is not merely to show that a gap can be made large, but to characterize the state variables that make posterior integration economically relevant.

\subsection{Empirical option panel}

The raw empirical sample contains every downloaded WTI APO history rather than a hand-selected set of strikes. The unit of observation is option $i$ on trading date $t$, identified by option type, strike, expiry month, and symbol. The raw panel contains 213 unique contracts and 11,179 option-date observations. The sample is intentionally unbalanced because strike availability and liquidity decline at longer maturities.

Three data layers remain separate. The \emph{raw panel} preserves all parsed observations. The \emph{main empirical panel} applies transparent data-quality and pricing-availability rules. The \emph{representative-contract set} is used only for figures and economic discussion. The October-2026 baseline contains 310 retained contract-date observations across 12 eligible valuation dates. Five dates contain at least one option with positive reported volume, and six individual contract-date observations report positive daily volume themselves.

\subsection{Reconstructing first-nearby fixings and moneyness}

Each option-date observation is linked to the information set required by the APO payoff. The pipeline reconstructs: (i) realized first-nearby settlements for fixing dates already elapsed, (ii) the relevant CL futures contracts for remaining fixing dates, and (iii) the contemporaneous futures term structure used to initialize those remaining fixings. This mapping produces $\widehat A^Q_{t,T}$ in Equation~\eqref{eq:expected_average} and the observation-specific moneyness measure in Equation~\eqref{eq:moneyness}.

The baseline physical-measure inference uses Yahoo \texttt{CL=F} as a labelled continuous/front-month return proxy; it is not used as a contractual substitute for the first-nearby sequence that determines the APO payoff. A dedicated source-robustness experiment reconstructs the historical first-nearby return series contract by contract from official final CL settlements. Each trading date is assigned to the earliest contract whose last-trade date has not passed. Returns spanning a contract switch are excluded so that contango or backwardation is not treated as a one-day diffusion shock. Four isolated mapped settlement gaps are retained in the audit trail but not imputed; any one-day return depending on a missing price is also excluded. The resulting 2024--2026 reconstruction contains 632 usable daily returns and 33 mapped roll switches, compared with 676 usable Yahoo returns over the same inference window.

Valuation uses committed Barchart individual CL histories with explicit last-trade dates. In the Barchart histories used here, \texttt{Latest} is the CME settlement field; the source-field name is retained for provenance and is not interpreted as an intraday last trade.

\subsection{Rolling physical-measure uncertainty}

Let $\mathcal D_t$ denote the historical information set available at date $t$. For an expanding sample beginning in January 2024, the posterior
\begin{equation}
 p(\sigma\mid\mathcal D_t)
\end{equation}
is estimated without using future returns. Each option-date observation is then priced using posterior integration, posterior-mean volatility, the marginal posterior mode, and the historical MLE. This rolling construction prevents later returns from contaminating earlier valuation dates.

The principal pricing errors are
\begin{equation}
 e^M_{i,t}=P^{M}_{i,t}-P^{mkt}_{i,t},
\end{equation}
where $P^{mkt}_{i,t}$ is the Barchart end-of-day settlement field and $M$ indexes the pricing rule. Mean error, MAE, and RMSE are reported overall and within liquidity-restricted samples.

To separate information source from information recency, a pre-specified recent-history benchmark reuses the exact strict-forward holdouts and target-date information cutoffs. It evaluates annualized realized volatility from the most recent 63, 126, and 252 usable daily returns, plus an exponentially weighted estimate with a 63-business-day half-life. These choices are fixed before inspecting their pricing errors. Each scalar estimate is mapped through the same Curran APO pricing engine used for the prior-date smile predictions. The purpose is not to optimize a historical window ex post, but to ask whether a reasonable recent historical benchmark materially narrows the option-informed advantage.

A stronger historical forecast benchmark uses a Gaussian GARCH(1,1), a standard class in crude-oil volatility forecasting \cite{gildertsiaras2020}. For each target date, parameters are estimated only from usable returns dated strictly before that date,
\begin{equation}
 h_{t+1}=\omega+\alpha\varepsilon_t^2+\beta h_t,
 \qquad \alpha+\beta<1.
 \label{eq:garch11}
\end{equation}
The observed target-date return is then used only to update the end-of-day conditional variance state, not to re-estimate parameters. Expected daily variances are forecast over the remaining business-day horizon and accumulated to each unresolved fixing date. Those cumulative variances define the common-factor covariance of the remaining futures fixings. Finally, a single annualized $\sigma_{GARCH}$ is chosen so that the constant-volatility pricing representation matches the arithmetic-average variance implied by the GARCH forecast path. The resulting horizon-matched scalar is evaluated on the same 2,381 strict-forward holdouts.

A second scalar sensitivity uses Equation~\eqref{eq:sigma_rms}. Because every empirical run already stores the posterior mean and standard deviation of $\sigma$, the posterior-RMS comparator can be evaluated on the same holdouts without rerunning MCMC.

\subsection{APO-implied volatility and cross-sectional validation}

The initial September and October panels provide 402 contract-date observations for implied-volatility analysis. Contract-level implied volatilities are obtained by inverting Equation~\eqref{eq:apo_iv}. A date-level common $\sigma_Q$ is then calibrated by minimizing cross-sectional settlement squared error. Two deliberately harder checks avoid evaluating a contract on information extracted from that same contract: leave-one-contract-out calibration excludes the target option, and cross-option-type transfer estimates from calls and evaluates puts or vice versa.

These exercises remain cross-sectional. The primary predictive robustness design is therefore strictly forward in time. For each target date and expiry, the volatility smile is estimated using only APO implied-volatility observations from dates strictly earlier than the target. Two schemes are considered: a \emph{previous-day smile} using the latest earlier completed date and an \emph{expanding smile} using all earlier dates with exponential recency weights. Same-day implied volatility is retained only as an ex-post diagnostic and never enters the target-date fit.

A persistence benchmark asks whether the fitted previous-day smile adds information beyond simply carrying forward the same contract's latest strictly prior APO implied volatility. The carried volatility is repriced using the target-date futures curve, realized-fixing state, discount factor, and remaining fixing schedule. Thus the benchmark preserves a transparent current-state adjustment while excluding same-day APO information. A committed-data coverage audit requires at least six main-sample contracts on a date before an expiry is admitted to the smile experiment. Seven expiries pass that screen---September, October, and November 2026; March and September 2027; and March and September 2028---while June 2029 does not. The extended evaluation contains 2,381 option-date holdouts over 15 distinct target dates, with results reported both pooled and separately by expiry.

\subsection{Independent vanilla-WTI forward validation}\label{sec:external_lo_design}

The APO prior-date smile establishes that option-market volatility information is forward-relevant, but it does not by itself rule out a within-family explanation. The independent validation therefore constructs a risk-neutral state from standard monthly WTI futures options (LO) and uses it to predict October-2026 APO settlements. Raw vendor records are kept local; only reproducible query metadata, diagnostics, and derived results are versioned.

The acquisition window is August 24 through September 10, 2026, with CME trading reference date used for temporal ordering. The final strike range is 70--120 on the two CL futures contracts required by the October APO fixing schedule, CLX6 and CLZ6. Final official LO settlements are paired with official CL futures settlements and inverted using Equation~\eqref{eq:lo_iv}. The resulting panel contains 5,220 successful contract-date implied-volatility inversions over 13 reference dates. Of these settlements, 5,214 are flagged actual and six theoretical.

For each target APO date $t$, every LO observation used to fit the external surface satisfies $s<t$. The \emph{previous-day} surface uses only the latest completed LO reference date; the \emph{expanding} surface uses all earlier reference dates with a five-calendar-day exponential half-life. The primary specification evaluates the fitted surface separately for CLX6 and CLZ6 at each APO strike and combines the component volatilities with the fixing-count RMS rule in Equation~\eqref{eq:lo_effective_sigma}. A pre-specified sensitivity instead chooses the scalar volatility that matches the variance of the unresolved arithmetic-average component under the maintained common-factor representation, using target-date futures levels and fixing times. Same-day APO prices are never used to estimate the external state.

The final October comparison is matched at the contract-date level. Across the 280 observations shared by the historical baseline, the two external LO specifications, and the two prior-date APO smiles, 253 contract-dates lie inside the observed prior LO moneyness support for both external surface specifications. This 90.4\% no-clipping subset spans all 10 matched valuation dates and is the primary external-validation sample. Error differences are resampled by valuation-date cluster with 100,000 bootstrap replications so that contracts sharing one market date are not treated as independent observations.

\subsection{Heavy-tail robustness under the physical measure}

To test whether the Gaussian historical likelihood drives the pricing hierarchy, the physical-measure innovation is replaced by a standardized Student-$t$ variable,
\begin{equation}
 R_i=
 \left(\mu-\frac12\sigma^2\right)\Delta t
 +\sigma\sqrt{\Delta t}\,\varepsilon_i,
 \qquad
 \operatorname{Var}(\varepsilon_i)=1,
\end{equation}
with degrees of freedom inferred jointly with $\mu$ and $\sigma$. The $Q$ pricing dynamics, realized fixings, futures curve, and pricing map are otherwise unchanged. The production robustness run uses eight chains of 100,000 iterations per unique valuation date and covers 402 option-date observations across the September and October expiries.

\subsection{Numerical accuracy and posterior calibration}

Small values of $\Delta_{PI,PM}$ are scientifically useful only if they are larger than pricing noise. Eighteen empirically difficult contracts are therefore repriced over a dense volatility grid with independent high-precision pseudo-random Monte Carlo runs. Eight common difficult targets are also evaluated with independently scrambled Sobol sequences, and Curran conditioning is evaluated on the same volatility grids. The budgets are chosen so that direct Monte Carlo uncertainty in the PI--PM gap is typically of order $10^{-6}$, well below the observed empirical gaps near $10^{-3}$.

Posterior implementation is checked separately by simulation-based calibration. Two hundred thousand datasets are drawn from the prior predictive distribution, with 50,000 replications at each of $n=21,63,252,1260$. Calibration is assessed through posterior-CDF ranks and 50\%, 80\%, and 95\% credible-interval coverage.

\subsection{Liquidity dependence and clustered uncertainty}

Contracts observed on the same valuation date share the futures curve, historical posterior, and market shock, so treating option rows as independent would overstate precision. Error comparisons are therefore bootstrapped by valuation-date cluster rather than by individual contract. The production robustness analysis uses 100,000 cluster resamples for the full sample, the positive-volume subset, and open-interest thresholds of 1, 10, 100, and 500 contracts. The bootstrap reports percentile intervals and the probability that a candidate specification improves MAE or RMSE relative to the historical-volatility posterior-integrated baseline.
